% ============================================================
%  Informe Técnico — Práctica Final Integradora
%  Análisis y Diseño de Algoritmos · Universidad EAFIT · 2026-1
% ============================================================
\documentclass[11pt, a4paper]{article}

\usepackage[spanish]{babel}
\usepackage[utf8]{inputenc}
\usepackage[T1]{fontenc}
\usepackage{microtype}

\usepackage[
  a4paper,
  top=2cm, bottom=2cm,
  left=2.5cm, right=2cm,
  headheight=14pt
]{geometry}

\usepackage{fancyhdr}
\pagestyle{fancy}
\fancyhf{}
\fancyhead[L]{\small\textit{Práctica Final Integradora — ADA 2026-1}}
\fancyhead[R]{\small\textit{Universidad EAFIT}}
\fancyfoot[C]{\small\thepage}
\renewcommand{\headrulewidth}{0.4pt}

\usepackage{amsmath, amssymb}
\usepackage{booktabs, tabularx, array}
\renewcommand{\arraystretch}{1.18}
\usepackage{enumitem}
\setlist{topsep=3pt, itemsep=2pt, parsep=0pt}

\usepackage{listings}
\lstset{
  basicstyle=\footnotesize\ttfamily,
  frame=single,
  numbers=left,
  numberstyle=\tiny\color{gray},
  numbersep=5pt,
  breaklines=true,
  captionpos=b,
  xleftmargin=8pt,
  aboveskip=5pt,
  belowskip=5pt,
  showstringspaces=false,
}
\usepackage{xcolor}
\usepackage{caption}
\captionsetup{font=small, labelfont=bf, skip=4pt}
\usepackage[colorlinks=true, linkcolor=black, urlcolor=black]{hyperref}

\usepackage{tikz}
\usetikzlibrary{shapes.geometric, arrows.meta, positioning, fit, calc}

\tikzset{
  bloque/.style={rectangle, rounded corners=4pt, draw=black, fill=gray!10,
                 text width=3.8cm, align=center, minimum height=0.8cm,
                 font=\small},
  io/.style={trapezium, trapezium left angle=70, trapezium right angle=110,
             draw=black, fill=blue!8, text width=3.2cm, align=center,
             minimum height=0.7cm, font=\small},
  flecha/.style={-{Stealth[length=5pt]}, thick},
}

\setlength{\parindent}{0pt}
\setlength{\parskip}{5pt}

% ============================================================
\begin{document}
\thispagestyle{empty}

\begin{center}
  \vspace*{1cm}
  {\large\bfseries Universidad EAFIT}\\[3pt]
  {\normalsize Departamento de Ciencias de la Computación}\\[2pt]
  {\normalsize Análisis y Diseño de Algoritmos · 2026-1}

  \vspace{1.8cm}
  \rule{\linewidth}{0.6pt}\\[8pt]
  {\LARGE\bfseries Práctica Final Integradora}\\[6pt]
  {\large Informe Técnico}\\[8pt]
  \rule{\linewidth}{0.6pt}

  \vspace{0.6cm}
  {\normalsize\itshape Optimización de rutas y planificación de recursos en redes de telecomunicaciones}

  \vspace{1.5cm}
  \begin{tabular}{rl}
    \textbf{Integrante 1:} & Nathalia Valentina Cardoza Azuaje \\[3pt]
    \textbf{Integrante 2:} & Andres Felipe Velez Alvarez \\[3pt]
    \textbf{Integrante 3:} & Sebastian Salazar Henao
  \end{tabular}

  \vspace{1cm}
  {\small Repositorio: \texttt{ADA\_PF\_Cardoza\_Salazar\_Velez}}\\[3pt]
  {\small 8 de mayo de 2026}
\end{center}

\newpage
\tableofcontents
\newpage

% ============================================================
\section{Introducción}
% ============================================================

El proyecto aplica tres paradigmas algorítmicos al dataset \textit{Telco Customer Churn}
(Kaggle, 7\,043 clientes de una operadora en California) para gestionar solicitudes de
servicio en una red de telecomunicaciones simulada. Cada módulo resuelve un subproblema
distinto, y la elección de paradigma en cada caso no es arbitraria: responde a propiedades
estructurales del problema que hacen que un enfoque sea correcto y los demás inadecuados.

\begin{itemize}
  \item \textbf{Módulo A --- Divide y Vencerás:} Ordenar y consultar solicitudes por
        antigüedad (\texttt{tenure}). MergeSort es correcto aquí porque el problema exhibe
        \textit{subestructura recursiva}: ordenar $n$ elementos equivale a ordenar dos
        mitades independientes y fusionarlas sin necesidad de reexaminar decisiones pasadas.
        Esta descomposición permite alcanzar $\Theta(n\log n)$, que es la cota inferior
        para cualquier algoritmo de ordenamiento por comparación, con estabilidad garantizada.

  \item \textbf{Módulo B --- Codicioso (Kruskal):} Diseñar la red de mínimo costo entre
        20 grupos de clientes. El MST posee la \textit{propiedad de elección codiciosa}:
        en cada paso, agregar la arista de menor peso que no forme ciclo conduce al óptimo
        global (Lema del Ciclo / Lema del Corte). Las decisiones pasadas no afectan la
        validez de las futuras; una vez unidas dos componentes, esa decisión nunca se revierte.

  \item \textbf{Módulo C --- Programación Dinámica:} Maximizar el ingreso del ancho de
        banda asignado a las 50 solicitudes más prioritarias. La Mochila 0-1 es NP-difícil
        y no admite elección codiciosa exacta: el ratio $v_i/w_i$ no garantiza el óptimo
        global porque incluir un ítem de alta densidad puede bloquear combinaciones más
        valiosas. Sin embargo, el problema tiene \textit{subestructura óptima} y
        \textit{subproblemas solapados}, las dos propiedades que justifican la PD y
        permiten obtener la solución exacta en tiempo pseudopolinomial $\Theta(nW)$.
\end{itemize}

% ============================================================
\section{Descripción del Dataset}
% ============================================================

El archivo \texttt{WA\_Fn-UseC\_-Telco-Customer-Churn.csv} contiene 7\,043 registros y
21 columnas. Se extraen cinco campos: \texttt{customerID} (identificador único),
\texttt{tenure} (antigüedad en meses, usado como prioridad), \texttt{MonthlyCharges}
(cargo mensual en USD), \texttt{TotalCharges} (cargo acumulado, peso en la mochila) y
\texttt{Churn} (estado activo/en riesgo).

\textbf{Estadísticas globales:} 7\,043 registros; 5\,174 activos (\texttt{Churn = No},
73.5\,\%); 1\,869 en riesgo (26.5\,\%). \texttt{tenure} entre 0 y 72 meses, media 32.4,
mediana 29. \texttt{MonthlyCharges} promedio 64.76 USD (rango 18.25--118.75).
\texttt{TotalCharges} promedio 2\,283 USD.

\textbf{Manejo de nulos:} 11 registros presentan \texttt{TotalCharges} como cadena vacía,
todos con \texttt{tenure = 0} (clientes sin historial de facturación acumulada). El parser
recorta espacios y caracteres \texttt{\textbackslash r} antes de la conversión numérica
y asigna $0.0$ a estos casos, preservándolos en lugar de descartarlos. Esta decisión es
conservadora: eliminarlos habría sesgado el conteo de solicitudes de \texttt{tenure = 0},
que son válidas aunque no hayan generado cargo total aún.

\textbf{Construcción del grafo (Módulo B):} Se forman 20 grupos asignando cada registro $i$
al grupo $i \bmod 20$, distribuyendo los 7\,043 clientes de forma aproximadamente uniforme
($\sim$352 por grupo). El grafo completo $K_{20}$ tiene $\binom{20}{2} = 190$ aristas con
peso $\lfloor \overline{MC}_u + \overline{MC}_v \rfloor$, modelando el costo de instalación
del enlace entre las zonas de servicio $u$ y $v$.

% ============================================================
\section{Módulo A --- Divide y Vencerás}
% ============================================================

\subsection{MergeSort (orden descendente por \texttt{tenure})}

\begin{lstlisting}[language=C++, caption=MergeSort y Merge]
MergeSort(arr, izq, der):
    si izq >= der: retornar                    // caso base: un elemento
    medio = (izq + der) / 2
    MergeSort(arr, izq, medio)                 // ordenar mitad izquierda
    MergeSort(arr, medio+1, der)               // ordenar mitad derecha
    Merge(arr, izq, medio, der)                // fusionar en orden descendente

Merge(arr, izq, medio, der):
    L = arr[izq..medio];  R = arr[medio+1..der]
    i = 0;  j = 0;  k = izq
    mientras i < |L| y j < |R|:
        si L[i].tenure >= R[j].tenure:         // >= preserva estabilidad
            arr[k++] = L[i++]
        si no:
            arr[k++] = R[j++]
    copiar restos de L en arr[k..]
    copiar restos de R en arr[k..]
\end{lstlisting}

\textbf{Análisis de complejidad:} La recurrencia $T(n)=2T(n/2)+\Theta(n)$, con $a=2$,
$b=2$ y $f(n)=\Theta(n)=\Theta(n^{\log_2 2})$, cae en el Caso~2 del Teorema Maestro:
$T(n)=\Theta(n\log n)$. Esta cota es \textbf{asintóticamente óptima}: cualquier
algorithmo de ordenamiento por comparación requiere $\Omega(n\log n)$ comparaciones en
el peor caso (demostrado por el argumento del árbol de decisión de altura mínima
$\lceil\log_2(n!)\rceil\approx n\log_2 n - n/\ln 2$). El costo de espacio auxiliar es
$\Theta(n)$ por los arreglos temporales de la fase \textsc{Merge}.

\textbf{Por qué MergeSort y no QuickSort:}
\begin{itemize}
  \item \textbf{Peor caso garantizado:} MergeSort es $\Theta(n\log n)$ en todo caso.
        QuickSort sin aleatorización alcanza $\Theta(n^2)$ sobre un arreglo ya ordenado
        (el pivote siempre queda en un extremo, produciendo particiones de tamaño $0$ y
        $n-1$), que es precisamente la situación de re-ejecución sobre el mismo CSV.
  \item \textbf{Estabilidad:} El comparador $\geq$ preserva el orden relativo de registros
        con igual \texttt{tenure} (hay 353 clientes con \texttt{tenure=1}). QuickSort
        básico no es estable y podría alterar el orden de duplicados, complicando la
        reproducibilidad de los resultados entre ejecuciones.
  \item \textbf{Determinismo:} MergeSort siempre produce exactamente el mismo resultado
        independientemente del orden inicial; QuickSort aleatorizado varía el orden de
        los duplicados entre ejecuciones.
\end{itemize}

\subsection{Búsqueda Binaria Recursiva}

\textbf{Especificación (corrección del enunciado):} La función busca un registro cuyo
\texttt{tenure} sea \textbf{exactamente igual} a $k$. Si existen varios registros con
el mismo valor, es válido retornar cualquiera de ellos. Si no existe ninguno, se reporta
\textit{no encontrado} (retorna $-1$). La condición $\geq k$ del enunciado original era
ambigua: en un arreglo descendente siempre retornaría el primer elemento para cualquier $k$,
haciendo el ejercicio trivial. Los cinco valores de consulta $\{72,60,45,30,12\}$ están
confirmados como existentes en el dataset, por lo que ninguna consulta retorna $-1$.

\begin{lstlisting}[caption=Búsqueda binaria -- coincidencia exacta en arreglo descendente]
BinSearch(arr, izq, der, k):
    // Pre: arr[izq..der] ordenado DESCENDENTEMENTE por tenure
    si izq > der: retornar -1              // intervalo vacio: no encontrado
    medio = (izq + der) / 2
    si arr[medio].tenure == k:             // coincidencia exacta
        retornar medio
    si arr[medio].tenure > k:             // k esta en valores menores (mitad derecha)
        retornar BinSearch(arr, medio+1, der, k)
    retornar BinSearch(arr, izq, medio-1, k)   // k en valores mayores (mitad izquierda)
\end{lstlisting}

\textbf{Análisis de complejidad:} La recurrencia $T(n)=T(n/2)+\Theta(1)$ se resuelve
por el Teorema Maestro (Caso 2, $f(n)=\Theta(1)=\Theta(n^{\log_2 1})$), dando
$T(n)=\Theta(\log n)$. Cada consulta sobre $n=7\,043$ registros requiere a lo sumo
$\lceil\log_2 7043\rceil = 13$ comparaciones. Si existen duplicados, la función retorna
cualquier ocurrencia válida (comportamiento correcto según la corrección del enunciado).

\subsection{Resultados empíricos}

El ratio teórico para MergeSort entre dos tamaños $n_1$ y $n_2$ es
$\dfrac{n_2\log_2 n_2}{n_1\log_2 n_1}$. Para $n_1=1\,000\to n_2=3\,500$:
$\frac{3500\times 11.77}{1000\times 9.97}=\frac{41\,195}{9\,970}\approx4.13\times$. Para
$n_2=3\,500\to n_3=7\,043$:
$\frac{7043\times 12.78}{3500\times 11.77}=\frac{90\,009}{41\,195}\approx2.18\times$.

\begin{table}[h!]
\centering\small
\begin{tabularx}{\textwidth}{
  >{\centering\arraybackslash}p{1.3cm}
  >{\centering\arraybackslash}p{2.1cm}
  >{\centering\arraybackslash}p{2.6cm}
  >{\centering\arraybackslash}p{2.0cm}
  >{\centering\arraybackslash}p{2.0cm}
  >{\centering\arraybackslash}X
}
\toprule
\textbf{$n$} & \textbf{MergeSort} & \textbf{Búsq.\ bin.\ (prom.\ $10^6$)} &
\textbf{Ratio emp.} & \textbf{Ratio teór.} & \textbf{Error rel.} \\
\midrule
1\,000 & 136\,$\mu$s   & 0.0110\,$\mu$s & --- (base) & --- & --- \\
3\,500 & 620\,$\mu$s   & 0.0085\,$\mu$s & 4.56$\times$ & 4.13$\times$ & 9.4\,\% \\
7\,043 & 1\,209\,$\mu$s & 0.0075\,$\mu$s & 1.95$\times$ & 2.18$\times$ & 10.6\,\% \\
\bottomrule
\end{tabularx}
\caption{Tiempos empíricos de MergeSort y Búsqueda Binaria ($10^6$ ejecuciones promedio).
Ratios respecto al tamaño de muestra anterior.}
\label{tab:tiempos}
\end{table}

Los ratios empíricos son consistentes con el crecimiento $O(n\log n)$: al pasar de
$1\,000$ a $3\,500$ el tiempo se multiplica por $4.56\times$ (teórico $4.13\times$) y
de $3\,500$ a $7\,043$ por $1.95\times$ (teórico $2.18\times$). Las diferencias respecto
al valor teórico puro se deben a efectos de caché y al overhead de \texttt{std::chrono},
típicos en mediciones de bajo nivel; la tendencia es coherente con $\Theta(n\log n)$.
Las 5 consultas (coincidencia exacta) retornan:
$k{=}72\to$\texttt{6904-JLBGY}, $k{=}60\to$\texttt{7426-WEIJX},
$k{=}45\to$\texttt{7925-PNRGI}, $k{=}30\to$\texttt{2684-EIWEO},
$k{=}12\to$\texttt{7450-NWRTR}.

% ============================================================
\section{Módulo B --- Algoritmo Codicioso (Kruskal)}
% ============================================================

\subsection{Algoritmo y Union-Find}

\begin{lstlisting}[caption=Kruskal con Union-Find por rango y compresión de camino]
Kruskal(G):
    ordenar aristas E por peso ascendente              // O(E log E)
    UF = UnionFind(G.numNodos)
    para cada arista (u, v, w) en orden ascendente:
        si NOT UF.conectados(u, v):                    // sin ciclo
            UF.unir(u, v)
            mst.agregar((u, v, w))
            si |mst| == numNodos - 1: romper
    retornar mst

find(x):                                               // compresion de camino
    si padre[x] != x: padre[x] = find(padre[x])
    retornar padre[x]

unir(a, b):                                            // union por rango
    ra = find(a);  rb = find(b)
    si rango[ra] < rango[rb]: padre[ra] = rb
    si no si rango[ra] > rango[rb]: padre[rb] = ra
    si no: padre[rb] = ra;  rango[ra]++
\end{lstlisting}

\textbf{Análisis de complejidad:} El paso dominante es el ordenamiento de las $E$
aristas: $O(E\log E)$. Cada operación de Union-Find (\textsc{Find} y \textsc{Union})
con compresión de camino y unión por rango tiene costo amortizado prácticamente
constante, de modo que las $O(E)$ operaciones de Union-Find no alteran el orden
asintótico. La complejidad total es $O(E\log E)$, que para $E\leq V^2$ equivale a
$O(E\log V)$. Para $V=20$ y $E=190$: el ordenamiento requiere $\approx190\times7.6\approx
1\,450$ comparaciones, y el MST queda determinado tras exactamente 19 inclusiones.

\subsection{Justificación: Lema del Ciclo aplicado al grafo generado}

\textbf{Lema del Ciclo:} Para cualquier ciclo $C$ en $G$, la arista de mayor peso en
$C$ no pertenece a ningún MST de $G$.

\textit{Demostración por contradicción.} Sea $e_{\max}=(u,v,w_{\max})$ la arista de
mayor peso en el ciclo $C$, y supóngase que $e_{\max}\in T$ para algún MST $T$.
Al eliminar $e_{\max}$ de $T$, éste se divide en dos componentes $S$ y $V\setminus S$.
El ciclo $C$ conecta $u\in S$ con $v\in V\setminus S$, por lo que debe existir al
menos otra arista $e'\in C\setminus\{e_{\max}\}$ que cruce el corte $(S,V\setminus S)$,
con $w(e')< w_{\max}$ (los pesos de nuestro grafo son casi todos distintos).
Entonces $T'=T-e_{\max}+e'$ es un árbol de expansión con
$w(T')=w(T)-(w_{\max}-w(e'))<w(T)$, contradiciendo que $T$ es mínimo. $\square$

\textbf{Ilustración con arista concreta.} Considérese el ciclo $C=(0,1,12,0)$
con aristas $(0{-}1,124)$, $(1{-}12,122)$, $(0{-}12,124)$. La arista de mayor peso
es $(0{-}12,124)$. El Lema del Ciclo predice que \emph{no pertenece a ningún MST}.
Efectivamente, Kruskal selecciona $(1{-}12,122)$ y $(0{-}1,124)$ para conectar estos
tres nodos; la arista $(0{-}12,124)$ es rechazada porque 0 y 12 ya están en la misma
componente Union-Find al momento de procesarla.

\textbf{Propiedad de elección codiciosa.} Kruskal garantiza el óptimo global
porque su invariante es el Lema del Corte: en cada iteración, la arista de menor
peso que cruza algún corte $(S, V\setminus S)$ pertenece a algún MST. Al descartar
aristas que formarían ciclo (aplicando el Lema del Ciclo), nunca rechaza una arista
que debería incluirse; al incluir la de menor peso disponible sin ciclo, siempre elige
una arista óptima del corte inducido. Ambas propiedades juntas garantizan que el
resultado es el MST global.

\subsection{Aristas del MST y peso total}

\begin{table}[h!]
\centering\small
\begin{tabular}{ccc|ccc}
\toprule
\textbf{Nodo $u$} & \textbf{Nodo $v$} & \textbf{Peso} &
\textbf{Nodo $u$} & \textbf{Nodo $v$} & \textbf{Peso} \\
\midrule
1 & 12 & 122 & 1 &  9 & 128 \\
12 & 16 & 123 & 1 & 10 & 128 \\
0 &  1 & 124 & 1 & 11 & 128 \\
2 & 12 & 127 & 1 & 13 & 128 \\
1 &  3 & 127 & 12 & 14 & 128 \\
1 &  4 & 127 & 1 & 17 & 129 \\
1 &  5 & 127 & 1 & 18 & 129 \\
1 &  6 & 127 & 1 & 19 & 129 \\
1 &  7 & 127 & 1 & 15 & 130 \\
1 &  8 & 127 & \multicolumn{2}{c}{\textbf{Total}} & \textbf{2\,388} \\
\bottomrule
\end{tabular}
\caption{19 aristas del MST en orden de selección por Kruskal (peso total 2\,388).}
\label{tab:mst}
\end{table}

La topología es estrellada en torno al nodo~1 (hub con $\overline{MC}_1$ más bajo
del grafo) y secundariamente al nodo~12. Toda arista $(1,v)$ tiene peso mínimo
entre todas las opciones para conectar el nodo $v$ a la componente ya formada.

\subsection{Verificación manual en subgrafo de 5 nodos}

Se calcula el MST del subgrafo inducido por $\{0,1,2,12,16\}$. Las aristas relevantes
entre estos nodos, ordenadas por peso, son:

\begin{center}\small
\begin{tabular}{ccc}
\toprule
\textbf{Arista} & \textbf{Peso} & \textbf{Decisión Kruskal} \\
\midrule
$(1,12)$ & 122 & Incluir — une comp.~$\{1\}$ y $\{12\}$ \\
$(12,16)$ & 123 & Incluir — une comp.~$\{1,12\}$ y $\{16\}$ \\
$(0,1)$ & 124 & Incluir — une comp.~$\{0\}$ y $\{1,12,16\}$ \\
$(0,12)$ & 124 & Rechazar — 0 y 12 ya conectados (ciclo) \\
$(2,12)$ & 127 & Incluir — une comp.~$\{2\}$ y $\{0,1,12,16\}$ \\
\bottomrule
\end{tabular}
\end{center}

Kruskal selecciona las 4 primeras aristas válidas con peso total
$122+123+124+127=496$. Cualquier otro árbol de expansión del subgrafo de 5 nodos
tiene peso $\geq 497$ (la siguiente arista disponible cuesta $\geq 124$). La
verificación confirma que Kruskal produce el MST correcto del subgrafo.

% ============================================================
\section{Módulo C --- Programación Dinámica (Mochila 0-1)}
% ============================================================

\subsection{Planteamiento y recurrencia}

Se toman las 50 solicitudes activas (\texttt{Churn = No}) de mayor \texttt{tenure}
del arreglo ya ordenado por el Módulo A. Cada solicitud $i\in\{1,\ldots,50\}$ tiene:
$w_i=\texttt{round(TotalCharges}_i\texttt{)}$ (peso) y
$v_i=\texttt{round(MonthlyCharges}_i\times 10)$ (valor). El objetivo es hallar el
subconjunto $S$ que maximice $\sum_{i\in S}v_i$ sujeto a $\sum_{i\in S}w_i\leq W$.

\textbf{Justificación de escenarios.} La entrega especifica $W=500$; sin embargo,
los pesos reales del top-50 de clientes activos por antigüedad superan las 500
unidades (clientes con 72 meses de servicio acumulan \texttt{TotalCharges} del
orden de miles de USD). Para observar el comportamiento completo del algoritmo se
evaluó también $W=5\,000$, que permite seleccionar ítems significativos y revelar
las diferencias entre PD y enfoque codicioso. Ambos escenarios se reportan.

La \textit{subestructura óptima} se establece así: si el ítem $n$ no está en la solución
óptima de $(n,W)$, dicha solución coincide con la óptima de $(n{-}1,W)$. Si sí está,
es el ítem $n$ más la solución óptima de $(n{-}1,W{-}w_n)$. La recurrencia:

\[
dp[i][w] =
\begin{cases}
  0 & \text{si } i=0 \text{ o } w=0 \\
  dp[i-1][w] & \text{si } w_i > w \\
  \max\!\bigl(dp[i-1][w],\ dp[i-1][w-w_i]+v_i\bigr) & \text{si } w_i \leq w
\end{cases}
\]

\subsection{Pseudocódigo de tabulación y backtracking}

\begin{lstlisting}[caption=Tabulación bottom-up de la Mochila 0-1]
Knapsack01(items, n, W):
    dp[0..n][0..W] = 0                          // tabla inicializada en cero
    para i = 1 hasta n:
        para w = 0 hasta W:
            si items[i].peso > w:
                dp[i][w] = dp[i-1][w]           // item i no cabe
            si no:
                dp[i][w] = max(
                    dp[i-1][w],                  // no tomar item i
                    dp[i-1][w - items[i].peso] + items[i].valor  // tomarlo
                )
    retornar dp[n][W]
\end{lstlisting}

\begin{lstlisting}[caption=Backtracking para recuperar los ítems seleccionados]
Backtrack(dp, items, n, W):
    seleccionados = [];  w = W
    para i = n hasta 1 (decreciente):
        si dp[i][w] != dp[i-1][w]:
            seleccionados.agregar(items[i])      // item i fue incluido
            w = w - items[i].peso                // descontar su peso
    retornar seleccionados
\end{lstlisting}

El backtracking recorre la tabla de $(n,W)$ hacia $(1,0)$ en $O(n+W)$: si
$dp[i][w]\neq dp[i-1][w]$, el ítem $i$ contribuyó al óptimo y se descuenta su peso.
La corrección está garantizada por la subestructura óptima: si el ítem $i$ pertenece
a la solución óptima de $(i,w)$, necesariamente $dp[i-1][w-w_i]+v_i>dp[i-1][w]$.

\subsection{Contraejemplo codicioso y discusión de pseudopolinomialidad}

Con $W=500$, todos los ítems del top-50 tienen $w_i>500$ (clientes de 72 meses acumulan
$\sim$6\,500 USD en \texttt{TotalCharges}), por lo que la solución óptima es vacía.
Este resultado no es un fallo: ilustra que la restricción de capacidad fue fijada por
debajo de la escala real del dataset; con instancias sintéticas calibradas,
$W$ habría sido elegido para que al menos algunos ítems quepan.

\begin{table}[h!]
\centering\small
\begin{tabularx}{\textwidth}{
  >{\raggedright\arraybackslash}p{3.5cm}
  >{\raggedright\arraybackslash}X
  >{\centering\arraybackslash}p{1.8cm}
  >{\centering\arraybackslash}p{1.5cm}
}
\toprule
\textbf{Enfoque} & \textbf{Solicitudes seleccionadas} & \textbf{Valor total} & \textbf{Óptimo?} \\
\midrule
$W=500$ --- PD & (ninguna) & 0 & Sí \\
$W=5000$ --- Codicioso ($v/w$) & 4312-KFRXN, 6734-PSBAW, 7606-BPHHN & 688 & No \\
$W=5000$ --- PD (Mochila 0-1) & 9286-BHDQG, 4312-KFRXN & \textbf{707} & Sí \\
\bottomrule
\end{tabularx}
\caption{Resultados comparativos de la Mochila para $W=500$ y $W=5\,000$.}
\end{table}

El codicioso por ratio $v/w$ selecciona 3 ítems (peso acumulado 4\,904, valor 688).
La PD descarta dos e incluye \texttt{9286-BHDQG} (mayor valor absoluto, menor ratio),
logrando 707. Este ejemplo muestra por qué la Mochila 0-1 no tiene propiedad de elección
codiciosa: la densidad $v_i/w_i$ no captura el efecto del peso residual sobre la capacidad
disponible para el resto de ítems.

\textbf{Análisis formal de complejidad.}
\begin{itemize}
  \item \textbf{Tiempo:} El doble bucle de tabulación itera exactamente
        $(n+1)\times(W+1)$ celdas, realizando $O(1)$ trabajo por celda (una
        comparación y un máximo). Por tanto:
        \[
          T(n,W)=\Theta(n\cdot W).
        \]
        Para la instancia del proyecto ($n=50$, $W=5\,000$): $50\times5\,001=250\,050$
        operaciones de comparación.
  \item \textbf{Espacio:} La tabla \texttt{dp[0..n][0..W]} ocupa
        $(n+1)(W+1)$ enteros: $\Theta(n\cdot W)$. Si sólo se necesita
        el valor óptimo (sin backtracking), basta con mantener dos filas alternas,
        reduciendo a $\Theta(W)$; en este proyecto se conserva la tabla completa para
        permitir la reconstrucción.
  \item \textbf{Pseudopolinomialidad:} El costo $\Theta(n\cdot W)$ es polinomial en
        el \emph{valor numérico} de $W$, pero $W$ se codifica en $\lceil\log_2 W\rceil$
        bits. Expresado en función del tamaño de la entrada en bits
        $s=\Theta(n\log W + n\log V_{\max})$, se tiene $W=2^{\Omega(s/n)}$, por lo que
        el algoritmo es \textbf{exponencial en $s$}. Por eso la Mochila 0-1 permanece
        NP-difícil a pesar de tener solución en $\Theta(nW)$: no existe algoritmo
        polinomial en $s$ salvo que P = NP.
\end{itemize}

% ============================================================
\section{Integración del Pipeline}
% ============================================================

\subsection{Diagrama de flujo del programa completo}

\vspace{4pt}
\begin{center}
\begin{tikzpicture}[node distance=0.55cm and 1.1cm]

  \node[io]          (csv)   {Leer CSV\\(7\,043 filas)};
  \node[bloque, below=of csv]  (parse) {parser.cpp:\\limpiar nulos,\\construir vector};
  \node[bloque, below=of parse](graph) {graph.cpp:\\20 grupos,\\190 aristas};
  \node[bloque, below=of graph](krus)  {kruskal.cpp:\\MST (19 aristas,\\peso 2\,388)};
  \node[io,    below=of krus] (mst)   {mst\_red.txt};

  \node[bloque, right=2.2cm of parse]  (msort)   {mergesort.cpp:\\ordenar por tenure\\descendente};
  \node[bloque, below=of msort]        (bsearch) {binary\_search.cpp:\\5 consultas\\$k\in\{72,60,45,30,12\}$};
  \node[io,    below=of bsearch]       (csv2)    {solicitudes\_ordenadas.csv\\busquedas\_A.txt};

  \node[bloque, right=2.2cm of msort]  (items)   {construirItemsMochila():\\top-50 activos\\del arreglo ordenado};
  \node[bloque, below=of items]        (kp)      {knapsack.cpp:\\PD $W=500$\\PD $W=5000$};
  \node[io,    below=of kp]            (bw)      {asignacion\_bw\_500.txt\\asignacion\_bw\_5000.txt};

  \draw[flecha] (csv)   -- (parse);
  \draw[flecha] (parse) -- (graph);
  \draw[flecha] (graph) -- (krus);
  \draw[flecha] (krus)  -- (mst);

  \draw[flecha] (parse.east) -- ++(0.4,0) |- (msort.west);
  \draw[flecha] (msort)  -- (bsearch);
  \draw[flecha] (bsearch)-- (csv2);

  \draw[flecha] (msort.east) -- ++(0.4,0) |- (items.west)
      node[midway, above, font=\tiny, xshift=20pt]{arreglo en memoria};
  \draw[flecha] (items)  -- (kp);
  \draw[flecha] (kp)     -- (bw);

  \node[above=0.15cm of csv,   font=\small\bfseries, text=gray]{Entrada};
  \node[above=0.15cm of msort, font=\small\bfseries, text=gray]{Módulo A};
  \node[above=0.15cm of items, font=\small\bfseries, text=gray]{Módulo C};
  \node[below=0.15cm of graph, xshift=-0.3cm, font=\small\bfseries, text=gray]{Módulo B};

\end{tikzpicture}
\end{center}
\vspace{4pt}

\subsection{Integración A $\rightarrow$ C en memoria}

El programa ejecuta los módulos en el orden A → B → C, aunque el flujo lógico de datos
es A → C. El vector \texttt{solicitudes}, una vez ordenado descendentemente por MergeSort,
permanece en memoria y es consumido directamente por \texttt{construirItemsMochila()} sin
releer el CSV. Esta función itera el vector una sola vez ($O(n)$) y selecciona los primeros
50 registros con \texttt{Churn = No}; dado el orden descendente por \texttt{tenure}, estos
son siempre los clientes activos de mayor antigüedad. La arquitectura en memoria elimina
lecturas de disco redundantes y garantiza coherencia: todos los módulos operan sobre el
mismo estado del dato.

% ============================================================
\section{Conclusiones}
% ============================================================

\begin{table}[h!]
\centering\small
\begin{tabularx}{\textwidth}{
  >{\raggedright\arraybackslash}p{2cm}
  >{\raggedright\arraybackslash}p{2.4cm}
  >{\raggedright\arraybackslash}p{2.8cm}
  >{\raggedright\arraybackslash}X
}
\toprule
\textbf{Paradigma} & \textbf{Algoritmo} & \textbf{Complejidad} & \textbf{Garantía de optimalidad} \\
\midrule
Divide y Vencerás & MergeSort + Búsqueda binaria & $\Theta(n\log n)$ / $\Theta(\log n)$ &
  Exacto y cota inferior óptima para ordenamiento por comparación. \\
Codicioso & Kruskal (MST) & $O(E\log E)$ &
  Exacto: Lema del Corte y Lema del Ciclo garantizan optimalidad global a partir de elecciones locales. \\
Prog. Dinámica & Mochila 0-1 & $\Theta(n\cdot W)$ &
  Exacto pero pseudopolinomial; NP-difícil en el tamaño real de la entrada. \\
\bottomrule
\end{tabularx}
\caption{Comparación de los tres paradigmas en términos de optimalidad y complejidad.}
\end{table}

\begin{itemize}
  \item \textbf{Elección de paradigma y estructura del problema:} Los tres módulos
        ilustran que la elección algorítmica correcta depende de propiedades estructurales
        del problema. El MST admite voracidad exacta porque sus decisiones son irrevocables
        e independientes. La Mochila no la admite porque incluir un ítem puede bloquear
        combinaciones más valiosas que solo se revelan al considerar el conjunto completo;
        la PD explora implícitamente todas las combinaciones mediante sus $n\times W$
        subproblemas solapados.

  \item \textbf{Datos reales vs. instancias sintéticas:} El escenario $W=500$ con valor
        óptimo 0 ilustra un fenómeno común con datasets reales: la escala de los pesos
        puede hacer que un algoritmo correcto produzca resultados triviales. Con instancias
        sintéticas calibradas, $W$ se fijaría para que al menos algunos ítems quepan. Por
        otro lado, la diferencia de 19 unidades de valor entre PD y codicioso ($W=5\,000$)
        es modesta en términos absolutos, pero suficiente para demostrar que los pesos reales
        no satisfacen la propiedad de intercambio. En datasets con pesos más heterogéneos,
        la brecha sería mayor y la diferencia entre paradigmas más evidente.

  \item \textbf{Codicioso exacto vs. heurístico:} Kruskal es codicioso exacto porque el
        MST cumple simultáneamente la propiedad de elección codiciosa (Lema del Corte) y
        la subestructura óptima. El codicioso por ratio $v/w$ en la Mochila es una heurística
        sin garantía formal: produce soluciones razonables pero potencialmente subóptimas
        en presencia de pesos heterogéneos, exactamente el escenario del dataset.

  \item \textbf{Pipeline en memoria:} La arquitectura sin re-lectura del CSV demuestra
        que el diseño de estructuras de datos compartidas puede mejorar tanto el rendimiento
        (evitar I/O redundante) como la corrección (garantizar que todos los módulos operen
        sobre el mismo estado consistente del dato).
\end{itemize}

% ============================================================
\section{Herramientas Utilizadas}
% ============================================================

C++17 (\texttt{g++ -std=c++17 -O2}), Visual Studio Code, Git/GitHub,
\LaTeX{} (pdflatex, paquetes \texttt{babel}, \texttt{booktabs}, \texttt{listings},
\texttt{tikz}). Se utilizó \textbf{OverLeaf} para la documentación y compilación del archivo \LaTeX{}

Se utilizó \textbf{Claude} (Anthropic, Claude Sonnet 4.6) y  Chat GPT (5.3) como asistente de IA generativa
para redacción y revisión de secciones del informe, enriquecimiento de pseudocódigos, análisis de los modulos,  construcción del diagrama de flujo en TikZ y mejora de código.

% ============================================================
\section{Referencias}
% ============================================================

\begin{enumerate}[label={[\arabic*]}]
  \item Blastchar. \textit{Telco Customer Churn}. Kaggle, 2018.
        \url{https://www.kaggle.com/datasets/blastchar/telco-customer-churn}
  \item Cormen, T.\,H., Leiserson, C.\,E., Rivest, R.\,L., \& Stein, C. (2022).
        \textit{Introduction to Algorithms} (4.ª ed.). MIT Press.
  \item Material del curso ADA --- Universidad EAFIT, 2026-1.
  \item Kleinberg, J., \& Tardos, É. (2006). \textit{Algorithm Design}. Pearson.
\end{enumerate}

% ============================================================
\section{Roles del Equipo}
% ============================================================

\begin{table}[h!]
\centering\small
\begin{tabularx}{\textwidth}{
  >{\raggedright\arraybackslash}p{4.2cm}
  >{\raggedright\arraybackslash}X
}
\toprule
\textbf{Integrante} & \textbf{Contribuciones específicas} \\
\midrule
\textbf{Nathalia Valentina Cardoza Azuaje} &
  \textbf{Módulo A.} \texttt{mergesort.cpp}, \texttt{binary\_search.cpp},
  \texttt{parser.cpp}, \texttt{output.cpp}: parseo del CSV, manejo de nulos,
  MergeSort estable, búsqueda binaria recursiva, medición de tiempos con
  promedio de $10^6$ ejecuciones. Redacción de las secciones del Módulo A. \\[4pt]
\textbf{Sebastian Salazar Henao} &
  \textbf{Módulo B.} \texttt{graph.cpp}, \texttt{graph.hpp},
  \texttt{kruskal.cpp}, \texttt{kruskal.hpp}: construcción del grafo $K_{20}$,
  Union-Find con compresión de camino y unión por rango, cálculo del MST,
  \texttt{guardarMST()}. Integración final en \texttt{main.cpp}, coordinación
  del repositorio Git. Redacción de las secciones del Módulo B. \\[4pt]
\textbf{Andres Felipe Velez Alvarez} &
  \textbf{Módulo C.} \texttt{knapsack.cpp}, \texttt{knapsack.hpp}:
  tabulación bottom-up, backtracking, contraejemplo codicioso, reportes
  \texttt{asignacion\_bw\_*.txt}. Redacción de las secciones del Módulo C
  y revisión de la integración A$\to$C. \\
\bottomrule
\end{tabularx}
\caption{Distribución de roles por módulo.}
\end{table}

\end{document}